Klasifikasi artikel berita ke taksonomi IPTC Media Topic dibutuhkan untuk menyaring jutaan artikel per hari. Namun pengklasifikasi IPTC multibahasa terbuka yang ada, berbasis XLM-RoBERTa-large, membutuhkan \val{eff.teacher_onnx_int8.latency_ms_p50_1c}--\val{eff.teacher.latency_ms_p50_1c}~ms per artikel pada satu \eng{core} CPU. Penelitian ini mengusulkan StaticKD, yaitu \eng{knowledge distillation} dari pengklasifikasi tersebut ke \eng{student} berbasis \eng{static embedding} multibahasa. Kepala linear \eng{student} dilipat sehingga seluruh model menjadi satu tabel berisi 17 logit per token. \eng{Teacher} melabeli \val{data.pool.docs} artikel tak berlabel dari \val{data.pool.langs} bahasa. Untuk evaluasi, kami membangun set uji gabungan berisi \val{ts.docs} artikel dari \val{ts.langs} bahasa dengan label yang independen dari \eng{teacher}: anotasi manusia serta rubrik dan kategori penerbit yang dipetakan ke pohon IPTC, disertai kontrol kontaminasi dan \eng{site-cluster bootstrap}. Pada set uji ini, StaticKD mencapai macro-F1 \val{statickd.test.macro} (\eng{teacher} \val{teacher.test.macro}; \val{statickd.test.pct_teacher}) dengan \eng{throughput} \val{eff.statickd.docs_per_sec_1c} artikel per detik per \eng{core}, sekitar \val{eff.statickd.speedup_onnx} kali lebih cepat dari \eng{teacher} ONNX int8, dalam model \val{final.size}~MB yang hanya membutuhkan NumPy. StaticKD juga lebih akurat daripada fastText dan TF-IDF + SVM yang didistilasi dengan data yang sama. Ablasi dengan lima \eng{seed} per varian menunjukkan bahwa keberhasilannya terutama berasal dari data distilasi multibahasa dan inisialisasi embedding yang selaras lintas bahasa, bukan dari \eng{soft label}. \eng{Ensemble} beberapa \eng{seed} dapat dilipat ke satu tabel tanpa biaya inferensi. Kelemahan utama ada pada bahasa yang tidak terwakili di data distilasi maupun pra-latih.
